\documentclass{article}

\usepackage{amsmath,amssymb}
\usepackage{xurl}
\usepackage[hidelinks]{hyperref}
\setlength{\emergencystretch}{2em}

% Shared commands for notes in docs/paper-gaps/.

\DeclareUnicodeCharacter{2080}{\ifmmode{}_{0}\else\textsubscript{0}\fi}
\DeclareUnicodeCharacter{2081}{\ifmmode{}_{1}\else\textsubscript{1}\fi}
\DeclareUnicodeCharacter{2082}{\ifmmode{}_{2}\else\textsubscript{2}\fi}
\DeclareUnicodeCharacter{2099}{\ifmmode{}_{n}\else\textsubscript{n}\fi}

\newcommand{\Lean}{\textsc{Lean}}
\ExplSyntaxOn
\NewDocumentCommand{\leanid}{m}{
  \ifmmode
    \textsf{\detokenize{#1}}
  \else
    \group_begin:
    \urlstyle{sf}
    \tl_set:Nn \l_tmpa_tl {#1}
    \tl_replace_all:Nnn \l_tmpa_tl {\_} {_}
    \exp_args:NV \path \l_tmpa_tl
    \group_end:
  \fi
}
\ExplSyntaxOff
\providecommand{\tr}{\operatorname{tr}}
\newcommand{\tnleanIssueBase}{https://github.com/LionSR/TNLean/issues/}
\newcommand{\tnleanPrBase}{https://github.com/LionSR/TNLean/pull/}
\newcommand{\tnleanDocBase}{https://sirui-lu.com/TNLean/docs/}
\newcommand{\ghissue}[1]{\href{\tnleanIssueBase#1}{GitHub issue \##1}}
\newcommand{\ghpr}[1]{\href{\tnleanPrBase#1}{PR \##1}}
\newcommand{\doclink}[1]{\href{\tnleanDocBase#1.html}{\path{#1}}}

% Dirac notation and the identity operator, as in blueprint/src/macros/common.tex.
% \providecommand keeps note-local definitions authoritative.
\usepackage{dsfont}
\providecommand{\ket}[1]{|#1\rangle}
\providecommand{\bra}[1]{\langle #1|}
\providecommand{\braket}[2]{\langle #1 | #2 \rangle}
\providecommand{\ketbra}[2]{|#1\rangle\!\langle #2|}
\providecommand{\Id}{\mathds{1}}
\providecommand{\one}{\mathds{1}}
\providecommand{\C}{\mathbb{C}}
\providecommand{\MN}[1]{M_{#1}(\C)}
\providecommand{\MD}{M_D(\C)}

% External referencing for paper sources.  The build helper writes sanitized
% auxiliary files under build/paper-aux/ and excludes duplicated source labels.
\usepackage{xr}

% Import a generated paper auxiliary file with a caller-supplied label prefix.
% The candidate paths support builds from docs/paper-gaps/, the repository root,
% and build/paper-aux/ itself.
\newcommand{\importpaperaux}[2]{%
  \IfFileExists{../../build/paper-aux/#2.aux}{%
    \externaldocument[#1]{../../build/paper-aux/#2}%
  }{%
    \IfFileExists{build/paper-aux/#2.aux}{%
      \externaldocument[#1]{build/paper-aux/#2}%
    }{%
      \IfFileExists{#2.aux}{%
        \externaldocument[#1]{#2}%
      }{}%
    }%
  }%
}

% General external reference macros
\newcommand{\paperref}[2]{\ref{#1-#2}}
\newcommand{\papereqref}[2]{(\ref{#1-#2})}

% CPSV16 source (1606.00608 / MPDO-22-12-17-2.tex) external document and shortcuts
\importpaperaux{cpsv16-}{1606.00608/MPDO-22-12-17-2-tnlean}
\newcommand{\cpsvref}[1]{\paperref{cpsv16}{#1}}
\newcommand{\cpsveqref}[1]{\papereqref{cpsv16}{#1}}

% Verdict marker: \gapnote{<kind>}{<status>}. Kind is the policy
% classification (clarification, local-correction, scope-restriction,
% unfaithful, false-source, open-gap); status is open, wip (elimination
% underway), resolved, or historical. Machine-read by the site index;
% prints a verdict line.
\newcommand{\gapnote}[2]{\par\noindent\textbf{Verdict:} #1 (#2).\par}


\title{The SAL to Eta-Local Structure Step in
Cirac--Perez-Garcia--Schuch--Verstraete 2017 Appendix C.2}
\author{}
\date{2026-07-11}

\begin{document}
\maketitle
\gapnote{open-gap}{resolved}

\section{Source statement}

Appendix C.2 of Cirac--P\'erez-Garc\'ia--Schuch--Verstraete proves that a
simple matrix product density operator (MPDO) satisfying the strong area law
(SAL) has the commuting nearest-neighbor product form
\begin{align}
    \sigma^{(N)}(\mathcal K)
        &\propto \prod_{n=1}^{N} B_{n,n+1},
    \label{eq:cpgsv17_mpdo_sal_zcl_eta_local_structure_source_product_form}
\end{align}
where each $B_{n,n+1}$ is positive semidefinite and all translated two-site
operators commute \cite[Appendix~C.2, Proposition~C.8]{Cirac2016MPDO_arXiv}.
In the local source file, this assertion is Proposition \texttt{3to4} at
\path{Papers/1606.00608/MPDO-22-12-17-2.tex:1571--1593}.

Let $U$ be the local unitary from Lemma \texttt{propSN} of
Ref.~\cite{Cirac2016MPDO_arXiv}, and define
\begin{align}
    \widetilde{\sigma}^{(N)}(\mathcal K)
        &=
        U^{\dagger\otimes N}
        \sigma^{(N)}(\mathcal K)
        U^{\otimes N}.
    \label{eq:cpgsv17_mpdo_sal_zcl_eta_local_structure_transformed_density}
\end{align}
The proof uses the sector decomposition obtained earlier in Appendix C.2:
\begin{align}
    \widetilde{\sigma}^{(N)}(\mathcal K)
        &=
        \bigoplus_{k_1,\ldots,k_N}
        \bigotimes_{n=1}^{N}\eta_{k_n,k_{n+1}},
    \label{eq:cpgsv17_mpdo_sal_zcl_eta_local_structure_sector_decomposition}
\end{align}
where $\eta_{k_n,k_{n+1}}\geq 0$ acts between the right subspace of site $n$
and the left subspace of site $n+1$
\cite[Appendix~C.2]{Cirac2016MPDO_arXiv}. The source defines
\begin{align}
    B_{n,n+1}
        &=
        \sum_{k_n,k_{n+1}}
        \operatorname{Id}_{H_{L,n}^{(k_n)}}
        \otimes \eta_{k_n,k_{n+1}}
        \otimes
        \operatorname{Id}_{H_{R,n+1}^{(k_{n+1})}},
    \label{eq:cpgsv17_mpdo_sal_zcl_eta_local_structure_bond_definition}
\end{align}
with the operator acting as the identity away from sites $n$ and $n+1$.
Equation~\ref{eq:cpgsv17_mpdo_sal_zcl_eta_local_structure_sector_decomposition}
then becomes the ordered product of the bonds in
Eq.~\ref{eq:cpgsv17_mpdo_sal_zcl_eta_local_structure_bond_definition}, and
their sector-local action gives commutativity of the translates.

\section{Current formalization}

The entropy implication at the beginning of Lemma C.2 of
Ref.~\cite{Cirac2016MPDO_arXiv} is formalized. For $N\geq 4$, SAL saturation
gives $I_1=I_2$, hence
\begin{align}
    S_{N-1}+S_1
        &=S_2+S_{N-2}.
    \label{eq:cpgsv17_mpdo_sal_zcl_eta_local_structure_entropy_equality}
\end{align}
This is equality in strong subadditivity for the normalized $(N-1)$-site
marginal partitioned into three consecutive regions of lengths $1$, $1$, and
$N-3$. The Hayashi characterization used in Appendix C.2 gives a direct-sum
tensor-product decomposition on the middle site
\cite{Cirac2016MPDO_arXiv}. Splitting the third region into its first site and
its remaining $N-4$ sites, then tracing out the latter, preserves the
decomposition. The first-three-site marginal therefore has the form asserted
in Lemma \texttt{Lsigma3} for every $N\geq 4$. The decomposition may be chosen
separately for each $N$.\footnote{The formal statements are
\leanid{MPOTensor.isSSAEquality_tripartite_of_isSAL} and
\leanid{MPOTensor.exists_etaStructure_reducedBlockState_three_of_isSAL};
the partial-trace preservation result is
\leanid{Entropy.exists_quantumMarkovDecomposition_rightMarginalAlong}.}

The desired commuting product form is recorded as an eta-local structure for
$\mathcal K$.\footnote{The formal structure is
\leanid{MPOTensor.EtaLocalStructureData}.} This structure contains a positive,
translation-invariant two-site bond and its realization of the finite-chain
MPO. The development derives the commuting-form property from this structure
and, after assuming zero correlation length (ZCL) separately, proves the
GSNNCH--ZCL consequence.

Earlier versions introduced a structure that combined unrelated local,
commuting-form, and doubled-index hypotheses. They then used the doubled-index
condition to derive a fusion formulation of the renormalization fixed-point
condition. Those declarations did not express the source's structural
fixed-point condition and have been removed. The remaining task was to
construct the operators $B_{n,n+1}$ from SAL.

\section{The SAL-to-positive-bond construction}

\subsection{Inverse-map contraction and sector factorization}

After deriving
Eq.~\ref{eq:cpgsv17_mpdo_sal_zcl_eta_local_structure_entropy_equality}, the
source applies the inverse tensor $\mathcal K^{-1}$ to the Hayashi
decomposition and proves Lemma \texttt{propSN}
\cite[Appendix~C.2]{Cirac2016MPDO_arXiv}. The algebraic part is formalized.
The concrete tensor $\mathcal K$ factors into sector tensors $r_k$ and $l_k$,
and the neighboring operators are initially defined, without a positivity
claim, by
\begin{align}
    \eta_{k,h}
        &=r_k l_h.
    \label{eq:cpgsv17_mpdo_sal_zcl_eta_local_structure_raw_neighbor_operator}
\end{align}
For every nonempty chain length, their cyclic contraction gives
\begin{align}
    \widetilde{\sigma}^{(N)}(\mathcal K)
        &=
        \bigoplus_{k_1,\ldots,k_N}
        \bigotimes_{n=1}^{N}\eta_{k_n,k_{n+1}}.
    \label{eq:cpgsv17_mpdo_sal_zcl_eta_local_structure_raw_cyclic_factorization}
\end{align}

For the four-site chain, the virtual contraction of the traced fourth site is
identified before applying the inverse tensor. Define
\begin{align}
    R_4
        &=
        \tr\!\left(\rho^{(4)}(\mathcal K)\right)^{-1}
        \sum_i \mathcal K^{i,i}.
    \label{eq:cpgsv17_mpdo_sal_zcl_eta_local_structure_normalized_four_site_tail}
\end{align}
The normalized three-site marginal satisfies
\begin{align}
    \sigma_3^{(4)}(\mathcal K)_{u,v}
        &=
        \tr\!\left(\mathcal K^{u,v}R_4\right).
    \label{eq:cpgsv17_mpdo_sal_zcl_eta_local_structure_four_site_marginal_closure}
\end{align}
If the four-site trace is nonzero, then $R_4\neq 0$, so there are indices
$\alpha,\beta$ for which
\begin{align}
    (R_4)_{\beta,\alpha}
        &\neq 0.
    \label{eq:cpgsv17_mpdo_sal_zcl_eta_local_structure_nonzero_tail_entry}
\end{align}
Thus the boundary matrix and nonzero entry selected before Equation
\texttt{formK} of Ref.~\cite{Cirac2016MPDO_arXiv} require neither SAL nor
injectivity as separate hypotheses.\footnote{The formal declarations are
\leanid{MPOTensor.normalizedFourSiteTail},
\leanid{MPOTensor.reducedBlockState_four_three_apply},
\leanid{MPOTensor.normalizedFourSiteTail_ne_zero}, and
\leanid{MPOTensor.exists_normalizedFourSiteTail_entry_ne_zero}.}

The double inverse-map contraction preceding the sector factorization is also
formalized. If $R$ denotes the virtual contraction of the remaining sites,
then
\begin{align}
    &\sum_{p_1,p_3}
        (\mathcal K^{-1})^{\alpha_1,\beta_1}_{p_1}
        (\mathcal K^{-1})^{\alpha_3,\beta_3}_{p_3}
        \tr\!\left(\mathcal K^{p_1}\mathcal K^{p_2}\mathcal K^{p_3}R\right)
    \notag\\
    &\qquad =
        \mathcal K^{p_2}_{\beta_1,\alpha_3}
        R_{\beta_3,\alpha_1}.
    \label{eq:cpgsv17_mpdo_sal_zcl_eta_local_structure_double_inverse_contraction}
\end{align}
This is the calculation at lines 1415--1438 of Appendix C.2 of
Ref.~\cite{Cirac2016MPDO_arXiv}.\footnote{The formal declarations are
\leanid{MPOTensor.inverseMapThreeSiteContraction} and
\leanid{MPOTensor.inverseMapThreeSiteContraction_eq}.}

The source display at line 1427 writes the remaining scalar as
$m_{\beta_3,\alpha_3}$. Direct contraction of the matrix units gives
\begin{align}
    \tr\!\left(E_{\alpha_1,\beta_1}
        \mathcal K^{p_2}
        E_{\alpha_3,\beta_3}R\right)
        &=
        \mathcal K^{p_2}_{\beta_1,\alpha_3}
        R_{\beta_3,\alpha_1}.
    \label{eq:cpgsv17_mpdo_sal_zcl_eta_local_structure_matrix_unit_contraction}
\end{align}
The scalar is therefore $m_{\beta_3,\alpha_1}$. The diagram defining
$m_{\beta,\alpha}$ and the subsequent choice of a nonzero pair
$(\alpha,\beta)$ agree with this correction. The repeated $\alpha_3$ is a
local index typo, not an additional mathematical hypothesis.

For a three-site closure against $R$, let
$A^{(k)}_{\alpha_1,\beta_1}$ and
$B^{(k)}_{\alpha_3,\beta_3}$ be the contractions of the inverse tensor
against $\rho_{A b_1}^{(k)}$ and $\rho_{b_2 c}^{(k)}$, respectively. In the
basis selected by the middle-site unitary, each diagonal sector satisfies
\begin{align}
    &R_{\beta_3,\alpha_1}
        \widetilde{\mathcal K}^{(k)}_{\beta_1,\alpha_3}
        ((l,r),(l',r'))
    \notag\\
    &\qquad =
        p_k
        A^{(k)}_{\alpha_1,\beta_1}(l,l')
        B^{(k)}_{\alpha_3,\beta_3}(r,r').
    \label{eq:cpgsv17_mpdo_sal_zcl_eta_local_structure_hayashi_sector_comparison}
\end{align}
The factor $p_k$ is explicit because the formal Hayashi decomposition stores
normalized left and right density operators. The proof commutes the outer
contractions through the unitary basis change and applies
Eq.~\ref{eq:cpgsv17_mpdo_sal_zcl_eta_local_structure_double_inverse_contraction}.
\footnote{The formal declarations are
\leanid{MPOTensor.IsThreeSiteClosure},
\leanid{MPOTensor.physicalSlice},
\leanid{MPOTensor.hayashiInverseLeft},
\leanid{MPOTensor.hayashiInverseRight}, and
\leanid{MPOTensor.inverseMap_hayashi_sector_comparison}.}

The three-site marginal of the normalized four-site chain is a three-site
closure against $R_4$. The off-diagonal Hayashi sectors of the transformed
physical slice vanish. Dividing
Eq.~\ref{eq:cpgsv17_mpdo_sal_zcl_eta_local_structure_hayashi_sector_comparison}
by a selected nonzero entry $R_{\beta_3,\alpha_1}$ produces tensors $l_k$ and
$r_k$ such that
\begin{align}
    U_B
    \mathcal K^{\,\cdot,\cdot}_{\beta_1,\alpha_3}
    U_B^\dagger
        &=
        \bigoplus_k
        (l_k)_{\beta_1}\otimes(r_k)_{\alpha_3}.
    \label{eq:cpgsv17_mpdo_sal_zcl_eta_local_structure_physical_slice_factorization}
\end{align}
This holds for all virtual indices $\beta_1,\alpha_3$. The required nonzero
entry exists whenever the four-site trace is nonzero. This is the
factorization at lines 1434--1438 of Appendix C.2 of
Ref.~\cite{Cirac2016MPDO_arXiv}.\footnote{The formal declarations are
\leanid{MPOTensor.isThreeSiteClosure_reducedBlockState},
\leanid{MPOTensor.inverseMap_hayashi_sector_offdiagonal},
\leanid{MPOTensor.sectorTensorL},
\leanid{MPOTensor.sectorTensorR},
\leanid{MPOTensor.physicalSlice_sector_factorization}, and
\leanid{MPOTensor.exists_physicalSlice_sector_factorization}.}

For a finite family of normal sectors, define the diagonal blocks used in
Equation \texttt{QkKjs} of Ref.~\cite{Cirac2016MPDO_arXiv} by
\begin{align}
    O_s^{(k)}(\beta,\alpha)
        &=
        \left[U_B\kappa^{(s)}_{\beta,\alpha}U_B^\dagger\right]_{k,k}.
    \label{eq:cpgsv17_mpdo_sal_zcl_eta_local_structure_normal_sector_block}
\end{align}
Every sector with nonzero Hayashi weight occurs in at least one normal sector.
Nonzero diagonal entries of the two trace-one Hayashi states give a nonzero
diagonal entry of the three-site state, and the normal-sector family closure
then has a nonzero summand. This argument does not require a closing-matrix
hypothesis.

The uniqueness cancellation is formalized with its essential hypothesis made
explicit. If every closing matrix $R_s$ is nonzero, the principal
BNT--Markov identity implies that two distinct normal sectors cannot both have
nonzero $k$-blocks. It therefore defines a unique label $j_k$ and projections
\begin{align}
    P_s
        &=
        \sum_{\substack{k:p_k\neq 0\\j_k=s}}Q_k,
    \label{eq:cpgsv17_mpdo_sal_zcl_eta_local_structure_bnt_projection}\\
    \sum_s P_s
        &=
        \sum_{k:p_k\neq 0}Q_k.
    \label{eq:cpgsv17_mpdo_sal_zcl_eta_local_structure_active_projection_resolution}
\end{align}
These projections are mutually orthogonal.\footnote{The relevant declarations
are
\leanid{MPOTensor.exists_bntMarkovBlockNonzero_of_probability_ne_zero},
\leanid{MPOTensor.existsUnique_bntMarkovBlockNonzero_of_probability_ne_zero},
\leanid{MPOTensor.bntSectorProjection}, and
\leanid{MPOTensor.sum_bntSectorProjection}.}

\subsection{Closing matrices and BNT sector selection}

Once the copy weights agree, the distinguished common weight is nonzero and
\begin{align}
    q_j
        &=
        \sum_q\mu_{j,q}^{N},
    \label{eq:cpgsv17_mpdo_sal_zcl_eta_local_structure_sector_weight_sum}\\
    q_j
        &=
        r_j\mu_j^N
        \neq 0.
    \label{eq:cpgsv17_mpdo_sal_zcl_eta_local_structure_sector_weight_nonzero}
\end{align}
These identities hold for every chain length $N$.\footnote{The relevant
declarations are
\leanid{MPSTensor.SectorDecomposition.commonWeight_ne_zero} and
\leanid{MPSTensor.SectorDecomposition.coeff_ne_zero_of_weight_copy_independent}.}

Simplicity implies that the physical-trace transfer
$\mathcal B_j=\sum_i A_j^{ii}$ is not nilpotent. Hence there is a positive tail
length $L$ and, for every representative, indices $\alpha_j,\beta_j$ such that
\begin{align}
    \widetilde m_j
        &=
        (\mathcal B_j^L)_{\alpha_j,\beta_j}
        \neq 0.
    \label{eq:cpgsv17_mpdo_sal_zcl_eta_local_structure_nonzero_trace_tail}
\end{align}
Taking the coefficient at total chain length $L+3$ gives
\begin{align}
    R_j
        &=
        q_j\mathcal B_j^L,
    \label{eq:cpgsv17_mpdo_sal_zcl_eta_local_structure_closing_matrix}\\
    (R_j)_{\alpha_j,\beta_j}
        &=
        q_j\widetilde m_j
        \neq 0.
    \label{eq:cpgsv17_mpdo_sal_zcl_eta_local_structure_closing_matrix_entry}
\end{align}
This is the construction at lines 1714--1718 of Appendix C.2 of
Ref.~\cite{Cirac2016MPDO_arXiv}.\footnote{The relevant declarations are
\leanid{MPSTensor.SectorDecomposition.exists_common_nonzero_physicalTraceTail_entry}
and
\leanid{MPSTensor.SectorDecomposition.exists_common_nonzero_threeSiteClosingMatrix_entry};
the matrix form is
\leanid{MPSTensor.SectorDecomposition.exists_common_threeSiteClosingMatrix_ne_zero}.}

The closing-matrix bridge includes the normalization. Equality of the complete
matrix-product-vector families transports the original tensor to its
horizontal BNT representation, and the normalized three-site reduced state
has closing matrices
\begin{align}
    \widehat R_j
        &=
        \tr\!\left(\rho^{(L+3)}\right)^{-1}
        q_j\mathcal B_j^L.
    \label{eq:cpgsv17_mpdo_sal_zcl_eta_local_structure_normalized_closing_matrix}
\end{align}
The four-site specialization used by the later argument takes $L=1$. SAL makes
the normalization scalar nonzero, so every $\widehat R_j$ is nonzero.
\footnote{The relevant declarations are
\leanid{MPSTensor.SectorDecomposition.normalizedThreeSiteClosingMatrix},
\leanid{MPOTensor.reducedBlockState_reindex_isThreeSiteFamilyClosure_of_sameMPV₂Pos},
and
\leanid{MPOTensor.reducedBlockState_four_threeSiteFamilyClosure_nonzero_closing}.}

Thus the generic cancellation theorem's nonzero-closing hypothesis is derived
for the family closure used by the source. The simultaneous specialization is
also formalized. After common blocking makes the BNT representatives
simultaneously one-letter spanning, biCF supplies the inverse and SAL supplies
the four-site Hayashi decomposition. Every active Markov sector then has a
unique BNT label, and the BNT-labelled sums of Hayashi projections are
orthogonal, mutually disjoint, and resolve the active Markov support.
\footnote{The relevant declaration is
\leanid{MPOTensor.exists_bntSectorProjectors_four_of_sameMPV₂Pos_isSAL}.}
This specialization needs no further closing-matrix hypothesis.

For every physical slice of the $i$-th normal representative,
\begin{align}
    P_sA_iP_t
        &=0
        \qquad\text{unless }s=t=i,
    \label{eq:cpgsv17_mpdo_sal_zcl_eta_local_structure_projector_slice_vanishing}\\
    P_iA_iP_i
        &=A_i.
    \label{eq:cpgsv17_mpdo_sal_zcl_eta_local_structure_matching_projector_slice}
\end{align}
The zero-weight Markov summands retained by the formal decomposition annihilate
all physical slices. Therefore
Eq.~\ref{eq:cpgsv17_mpdo_sal_zcl_eta_local_structure_matching_projector_slice}
does not require an ambient resolution of the identity. Applying the
projections sitewise to every nonempty periodic chain gives
\begin{align}
    P_s^{\otimes N}
    \sigma^{(N)}(K)
    P_s^{\otimes N}
        &=
        n_s\sigma^{(N)}(\mathcal K_s),
        \qquad N\geq 1,
    \label{eq:cpgsv17_mpdo_sal_zcl_eta_local_structure_sitewise_sector_compression}
\end{align}
where the common copy weight is absorbed into $\mathcal K_s$.\footnote{The
relevant declarations are
\leanid{MPOTensor.bntSectorProjection_mul_physicalSlice_mul_eq_zero},
\leanid{MPOTensor.sum_bntSectorProjection_mul_physicalSlice_mul_self},
\leanid{MPOTensor.changePhysicalBasis_bntSectorProjection_basis},
\leanid{MPOTensor.sitewise_bntSectorProjection_mul_mpo_mul_self}, and
\leanid{MPOTensor.exists_bntProjectorSelection_positiveLength_of_sameMPV₂Pos_isSAL}.}

The literal identity resolution in Equation \emph{Pis} of
Ref.~\cite{Cirac2016MPDO_arXiv} is recovered in the ambient physical space.
Choose a BNT label $s_0$ from a copy with unit-modulus weight. Assign every
zero-weight Markov summand to $s_0$, while retaining the unique source label on
each positive-weight summand. Define
\begin{align}
    \widehat P_s
        &=
        \sum_{k:\widehat j_k=s}Q_k.
    \label{eq:cpgsv17_mpdo_sal_zcl_eta_local_structure_completed_projection}
\end{align}
The completed projections satisfy
\begin{align}
    \sum_s\widehat P_s
        &=\Id,
    \label{eq:cpgsv17_mpdo_sal_zcl_eta_local_structure_completed_identity_resolution}\\
    \widehat P_s\widehat P_t
        &=0,
        \qquad s\neq t.
    \label{eq:cpgsv17_mpdo_sal_zcl_eta_local_structure_completed_projection_orthogonality}
\end{align}
Because every zero-weight $Q_k$ annihilates all physical slices, the
reassignment preserves
\begin{align}
    \widehat P_s K_i^{\beta\alpha}\widehat P_t
        &=0
        \qquad
        \text{whenever }s\neq t\text{ or }i\neq s.
    \label{eq:cpgsv17_mpdo_sal_zcl_eta_local_structure_completed_slice_vanishing}
\end{align}
The retained zero-weight summands account exactly for the difference between
the active-support resolution and Equation \emph{Pis} of
Ref.~\cite{Cirac2016MPDO_arXiv}. The construction and its SAL specialization
are
\leanid{MPOTensor.completedBntSectorProjection} and
\leanid{MPOTensor.exists_bntSeparatingProjectors_of_sameMPV₂Pos_isSAL_of_weight_copy_independent}.

The source's standing assumptions are recovered by
\leanid{MPOTensor.exists_bntSeparatingProjectors_of_horizontalCF_isSourceZCL_isSAL}.
This declaration applies the ZCL equation to the horizontal canonical-form
gauge identity, uses simplicity to exclude a zero physical-trace transfer, and
derives the equality of copy weights proved at lines 1646--1665 of Appendix
C.2 of Ref.~\cite{Cirac2016MPDO_arXiv}. It then invokes the copy-independent
specialization. The final theorem therefore assumes only the simple biCF, ZCL,
and SAL context of Appendix C.2, not copy independence separately.

The empty-word value $N=0$ is excluded because its closed virtual trace records
the bond dimension rather than a physical chain. Positivity and ZCL of every
selected absorbed BNT representative are formalized. Positivity follows by
compressing the positive full-chain operator and dividing by the positive
integer $n_s$. ZCL follows by restricting the full transfer equation to the
$s$-th canonical block and using nonnilpotence to exclude the zero transfer.
\footnote{The relevant declarations are
\leanid{MPOTensor.commonWeightAbsorbedBasisMPOTensor_isMPDO_of_projectorSelection},
\leanid{MPOTensor.commonWeightAbsorbedBasisMPOTensor_isMPDO_of_sameMPV₂Pos_isSAL},
and
\leanid{MPOTensor.commonWeightAbsorbedBasisMPOTensor_isSourceZCL}.}

The strict normalization required in the sector entropy identity is also
formalized. An MPDO tensor with source ZCL has positive trace at every
nonempty chain length because its normalized physical-trace transfer is a
nonzero idempotent. For the absorbed BNT sectors this proves
\begin{align}
    p_j^{(N)}
        &>0,
    \label{eq:cpgsv17_mpdo_sal_zcl_eta_local_structure_sector_probability_positive}
\end{align}
rather than merely $p_j\geq 0$.\footnote{The general declarations are
\leanid{MPOTensor.trace_mpo_ne_zero_of_isSourceZCL} and
\leanid{MPOTensor.trace_mpo_pos_of_isMPDO_isSourceZCL}; their specialization is
\leanid{MPOTensor.trace_mpo_commonWeightAbsorbedBasisMPOTensor_pos}.}

The weighted orthogonal-direct-sum entropy formula used at lines 1760--1770 of
Appendix C.2 of Ref.~\cite{Cirac2016MPDO_arXiv} is formalized as
\leanid{vonNeumannEntropy_blockDiagonal_smul}, together with the
pairwise-annihilating-support form
\leanid{vonNeumannEntropy_eq_sum_of_pairwise_annihilating_supports}.
The same sector probabilities occur in all four marginals, the Shannon terms
cancel, and strict positivity together with strong subadditivity forces
equality in each sector. Thus every absorbed BNT representative satisfies SAL
under the biCF hypothesis imposed at the beginning of Case II in the source.
\footnote{The principal declarations are
\leanid{MPOTensor.blockEntropy_eq_sum_bntSectorProbability},
\leanid{MPOTensor.blockEntropy_add_blockEntropy_eq_bntSector}, and
\leanid{MPOTensor.commonWeightAbsorbedBasisMPOTensor_isSAL_of_sameMPV₂Pos}.}
This completes the entropy argument at lines 1748--1781 of Appendix C.2 of
Ref.~\cite{Cirac2016MPDO_arXiv} under its stated biCF hypothesis.

That passage contains a local endpoint typo. Line 1771 writes
$m<\lfloor N/2\rfloor$, but Definition 4.6 of
Ref.~\cite{Cirac2016MPDO_arXiv} gives the equality chain through
$I_{\lfloor N/2\rfloor}$, and lines 1781--1783 compare $m=L$ with $m=L+1$.
When $L+1=\lfloor N/2\rfloor$, the comparison requires the endpoint. The
formal entropy equivalence therefore uses
\begin{align}
    1
        &\leq m
        \leq \lfloor N/2\rfloor.
    \label{eq:cpgsv17_mpdo_sal_zcl_eta_local_structure_entropy_endpoint_range}
\end{align}

\subsection{Neighboring operators and the Hayashi decomposition}

The neighboring operators are assembled from the sector tensors as
\begin{align}
    \eta_{k,h}
        &=
        \sum_{\gamma}
        (r_k)_{\gamma}\otimes(l_h)_{\gamma}.
    \label{eq:cpgsv17_mpdo_sal_zcl_eta_local_structure_sector_eta_definition}
\end{align}
The contraction of two neighboring factorized slices places $\eta_{k,h}$
between the outer sector tensors. Its trace is
\begin{align}
    T_{k,h}
        &=
        \tr(\eta_{k,h}),
    \label{eq:cpgsv17_mpdo_sal_zcl_eta_local_structure_trace_matrix_entry}\\
    T_{k,h}
        &=
        (r_k\mid l_h),
    \label{eq:cpgsv17_mpdo_sal_zcl_eta_local_structure_closed_sector_pairing}
\end{align}
where $(r_k\mid l_h)$ denotes the pairing of the closed sector tensors obtained
by tracing the physical legs.\footnote{The formal declarations are
\leanid{MPOTensor.sectorEta},
\leanid{MPOTensor.physicalSlice_neighboring_contraction},
\leanid{MPOTensor.closedSectorL},
\leanid{MPOTensor.closedSectorR},
\leanid{MPOTensor.trace_sectorEta}, and
\leanid{MPOTensor.ExplicitEtaOperators.ofSectorTensors}.}

At lines 1446--1450 of Appendix C.2, the source derives positivity of
$\eta_{k,h}$ from the all-length projected chain identity and chooses the
sector tensors so that each factor is positive semidefinite
\cite{Cirac2016MPDO_arXiv}. The all-length algebraic identity is formalized,
but by itself it does not choose coherent positive representatives for all
neighboring pairs. Positive semidefiniteness of the selected representatives
therefore initially remained a hypothesis of the explicit eta-data
construction.

The summary figure \path{MPDO_etahk=RhLk.png} at source line 1391 draws $r_h$
next to $l_k$ while naming the operator $\eta_{k,h}$. The later definition
\texttt{etarl}, the identity in
Eq.~\ref{eq:cpgsv17_mpdo_sal_zcl_eta_local_structure_closed_sector_pairing},
and the products $\eta_{k,l}\otimes\eta_{l,h}$ in Proposition \texttt{3to5}
of Ref.~\cite{Cirac2016MPDO_arXiv} consistently establish the convention
$\eta_{k,h}=r_k l_h$. The earlier figure contains an index-swap typo.

The Hayashi decomposition of the concrete three-site marginal is formalized
for every $N\geq 4$. Independent bijections on the first two tensor factors
preserve equality in strong subadditivity because they carry each marginal to
the corresponding reindexed marginal without changing its entropy. The
canonical bijection between a singleton block and one physical-site index
therefore transports SAL equality from flattened block coordinates to the
tripartite coordinates of the $(N-1)$-site state. The forward Hayashi
characterization gives a decomposition before the final $N-4$ sites of the
third region are traced out. Partial trace preserves the middle-site direct
sum, and the contiguous-block reduction identity gives
\begin{align}
    \sigma_3^{(N)}(\mathcal K)
        &=
        \tr_{4,\ldots,N-1}\!\left[\sigma_{N-1}^{(N)}(\mathcal K)\right].
    \label{eq:cpgsv17_mpdo_sal_zcl_eta_local_structure_three_site_reduction}
\end{align}
This proves Lemma \texttt{Lsigma3} of Ref.~\cite{Cirac2016MPDO_arXiv} at every
admissible chain length without asserting that the decomposition is common to
different values of $N$.\footnote{The formal declarations are
\leanid{isSSAEquality_submatrix_prodEquiv} and
\leanid{MPOTensor.exists_etaStructure_reducedBlockState_three_of_isSAL}.}

The inverse-map tensors also determine a source-minimal physical-sector
factorization. The nonzero dimensions of its left and right factors follow
from the trace-one Hayashi density matrices and are not separate assumptions.
For an injective tensor satisfying SAL, the four-site Hayashi decomposition
and a nonzero entry of the normalized tail produce this raw factorization
without ZCL or positivity hypotheses. For any such factorization, every
complete nonempty cyclic product of neighboring operators is positive
semidefinite because it is a diagonal sector compression of the corresponding
finite-chain operator.\footnote{The formal declarations are
\leanid{MPOTensor.inverseMapPhysicalSectorFactorization},
\leanid{MPOTensor.inverseMapPhysicalSectorFactorization_neighboringOperator_eq_sectorEta},
\leanid{MPOTensor.nonempty_physicalSectorFactorization_of_isSAL},
\leanid{MPOTensor.exists_physicalSectorFactorization_with_cyclicNeighboringProduct_posSemidef_of_isSAL},
\leanid{MPOTensor.etaOfSectorTensors},
\leanid{MPOTensor.sum_cyclic_sectorTensor_eq_prod_etaOfSectorTensors},
\leanid{MPOTensor.trace_sector_word_eq_prod_etaOfSectorTensors},
\leanid{MPOTensor.mpo_submatrix_sector_eq_cyclicEtaTensorProduct}, and
\leanid{MPOTensor.mpo_reindex_sectorChainEquiv_eq_blockDiagonal_cyclicEta}.
For the source-minimal physical-sector factorization, the corresponding direct
statements are
\leanid{MPOTensor.PhysicalSectorFactorization.mpo_submatrix_sector_eq_cyclicNeighboringProduct}
and
\leanid{MPOTensor.PhysicalSectorFactorization.mpo_reindex_sectorChainEquiv_eq_blockDiagonal}.}

The positive choice of the individual $\eta_{k,h}$ and primitivity of the
active-sector trace matrix are obtained from this factorization. Positivity of
complete cyclic products alone does not select coherent positive
representatives for all neighboring pairs. Recurrence, injective spanning, and
a directed-cut argument control the choices across the entire sector graph.

\subsection{Pairwise positivity and its limitations}

For an MPDO, the chain of the physically conjugated tensor is a congruence of
the original chain. Every cyclic tensor product of neighboring operators is
therefore a diagonal sector block of a positive operator:
\begin{align}
    0
        &\leq
        \left[Q_{k_1}\otimes\cdots\otimes Q_{k_N}\right]
        \widetilde{\sigma}
        \left[Q_{k_1}\otimes\cdots\otimes Q_{k_N}\right],
    \label{eq:cpgsv17_mpdo_sal_zcl_eta_local_structure_projected_chain_positive}\\
    \left[Q_{k_1}\otimes\cdots\otimes Q_{k_N}\right]
    \widetilde{\sigma}
    \left[Q_{k_1}\otimes\cdots\otimes Q_{k_N}\right]
        &=
        \bigotimes_{n=1}^{N}\eta_{k_n,k_{n+1}}.
    \label{eq:cpgsv17_mpdo_sal_zcl_eta_local_structure_projected_chain_tensor_product}
\end{align}
The length-one and length-two cases imply
\begin{align}
    \eta_{k,k}
        &\geq 0,
    \label{eq:cpgsv17_mpdo_sal_zcl_eta_local_structure_diagonal_eta_positive}\\
    \eta_{k,h}\otimes\eta_{h,k}
        &\geq 0.
    \label{eq:cpgsv17_mpdo_sal_zcl_eta_local_structure_two_cycle_positive}
\end{align}
A rescaling lemma for a nonzero positive semidefinite Kronecker product then
shows that, whenever both $\eta_{k,h}$ and $\eta_{h,k}$ are nonzero, there is
a scalar $c\neq 0$ such that
\begin{align}
    c\,\eta_{k,h}
        &\geq 0,
    \label{eq:cpgsv17_mpdo_sal_zcl_eta_local_structure_forward_pair_rescaling}\\
    c^{-1}\eta_{h,k}
        &\geq 0.
    \label{eq:cpgsv17_mpdo_sal_zcl_eta_local_structure_reverse_pair_rescaling}
\end{align}
Under the additional symmetric-support condition that $\eta_{k,h}=0$ implies
$\eta_{h,k}=0$, one scalar per unordered pair gives a positive semidefinite
family that agrees with the raw family on the diagonal and preserves every
two-site block $\eta_{k,h}\otimes\eta_{h,k}$.\footnote{The formal declarations
are
\leanid{Matrix.exists_smul_posSemidef_of_kronecker_posSemidef},
\leanid{MPOTensor.mpo_conjugatePhysical_eq},
\leanid{MPOTensor.cyclicEtaTensorProduct_posSemidef},
\leanid{MPOTensor.sectorEta_self_posSemidef},
\leanid{MPOTensor.sectorEta_kronecker_posSemidef},
\leanid{MPOTensor.exists_smul_posSemidef_sectorEta}, and
\leanid{MPOTensor.exists_explicitEtaOperators_sectorEta}.}

This pairwise construction has two limitations. Symmetric support is not a
source hypothesis, and the construction preserves only chain blocks of length
at most two. The source instead rescales the sector tensors, which preserves
all lengths simultaneously. The formal argument removes these limitations by
reparameterizing zero-weight sectors, proving recurrence on the remaining
support, and applying one vertex rephasing to the whole sector graph.

The obstruction already occurs when every neighboring bond space is
one-dimensional. Take sectors $0,1,2,3$ and define
\begin{align}
    \eta_{k,k}
        &=1,
        \qquad k\in\{0,1,2,3\},
    \label{eq:cpgsv17_mpdo_sal_zcl_eta_local_structure_phase_obstruction_diagonal}\\
    \eta_{0,1}
        &=\eta_{0,2}
        =\eta_{1,3}
        =1,
    \label{eq:cpgsv17_mpdo_sal_zcl_eta_local_structure_phase_obstruction_real_edges}\\
    \eta_{2,3}
        &=i,
    \label{eq:cpgsv17_mpdo_sal_zcl_eta_local_structure_phase_obstruction_imaginary_edge}
\end{align}
with every other off-diagonal entry zero. Every directed cyclic product is
zero or one and is therefore positive semidefinite. A sector rephasing that
preserves each same-sector product has the form
\begin{align}
    \eta_{k,h}
        &\longmapsto
        z_k^{-1}z_h\eta_{k,h},
        \qquad |z_k|=1.
    \label{eq:cpgsv17_mpdo_sal_zcl_eta_local_structure_sector_rephasing}
\end{align}
Positivity on the edges $0\to 1$, $0\to 2$, and $1\to 3$ forces
\begin{align}
    z_1
        &=z_0,
    \label{eq:cpgsv17_mpdo_sal_zcl_eta_local_structure_phase_constraint_one}\\
    z_2
        &=z_0,
    \label{eq:cpgsv17_mpdo_sal_zcl_eta_local_structure_phase_constraint_two}\\
    z_3
        &=z_1.
    \label{eq:cpgsv17_mpdo_sal_zcl_eta_local_structure_phase_constraint_three}
\end{align}
The transformed edge $2\to 3$ therefore retains phase $i$. No coherent choice
of sector phases makes all four nonzero neighboring operators positive.

A sufficient condition for coherent rephasing is recurrence of the nonzero
support: every nonzero directed edge lies on a nonzero directed cycle.
Positivity of all cyclic tensor products then shows that each edge is a phase
times a positive operator, forces trivial phase around every cycle, and permits
a coherent choice of vertex phases. The source does not state recurrence as a
hypothesis. The formal development derives it from injectivity after
reparameterizing zero-weight sectors, without using the later primitivity
argument.

For the formal notion of a Hayashi decomposition, recurrence can fail before
this reparameterization because zero sector weights are permitted. Consider
the scalar-bond tensor on a two-dimensional physical space given by
\begin{align}
    K^{00}
        &=1,
    \label{eq:cpgsv17_mpdo_sal_zcl_eta_local_structure_zero_weight_example_nonzero_slice}\\
    K^{01}
        &=K^{10}
        =K^{11}
        =0.
    \label{eq:cpgsv17_mpdo_sal_zcl_eta_local_structure_zero_weight_example_zero_slices}
\end{align}
This tensor is injective and generates the rank-one projector onto the
constant-zero word at every length. Decompose the middle physical space into
its two one-dimensional coordinate sectors, assign weights
\begin{align}
    p_0
        &=1,
    \label{eq:cpgsv17_mpdo_sal_zcl_eta_local_structure_zero_weight_example_active_weight}\\
    p_1
        &=0,
    \label{eq:cpgsv17_mpdo_sal_zcl_eta_local_structure_zero_weight_example_inactive_weight}
\end{align}
and choose normalized left and right states in both sectors supported on the
zero coordinate. For the scalar tail $R=1$, the inverse-map sector tensors
satisfy
\begin{align}
    r_0
        &=r_1
        =1,
    \label{eq:cpgsv17_mpdo_sal_zcl_eta_local_structure_zero_weight_example_right_tensors}\\
    l_0
        &=1,
    \label{eq:cpgsv17_mpdo_sal_zcl_eta_local_structure_zero_weight_example_active_left_tensor}\\
    l_1
        &=0,
    \label{eq:cpgsv17_mpdo_sal_zcl_eta_local_structure_zero_weight_example_inactive_left_tensor}\\
    \eta_{k,h}
        &=p_h.
    \label{eq:cpgsv17_mpdo_sal_zcl_eta_local_structure_zero_weight_example_neighbor_operator}
\end{align}
Thus $1\to 0$ is a nonzero support edge, but no edge enters sector $1$. This
edge is not recurrent. Consequently recurrence cannot be derived for an
arbitrary decomposition satisfying only $p_k\geq 0$.

The source-faithful correction retains the full physical direct sum because a
zero-weight summand may have nonzero left and right factor dimensions. In the
selected inverse-map factorization, $p_k=0$ already implies $l_k=0$. Replacing
$r_k$ by zero on exactly those sectors leaves every product
$l_k\otimes r_k$ unchanged and removes the spurious outgoing edges. Incoming
edges already vanish because the target left tensor is zero. This
reparameterization and its neighboring-operator formula are formalized by
\leanid{MPOTensor.PhysicalSectorFactorization.zeroRightTensorOnInactive},
\leanid{MPOTensor.sectorTensorL_eq_zero_of_weight_eq_zero},
\leanid{MPOTensor.zeroWeightReparameterizedInverseMapPhysicalSectorFactorization},
and
\leanid{MPOTensor.zeroWeightReparameterizedInverseMapPhysicalSectorFactorization_neighboringOperator}.

Every neighboring operator incident to an inactive sector then vanishes, while
active--active operators remain unchanged. On the active support, recurrence
follows from an algebraic triangle-closing argument once the virtual matrices
provided by all sector entries span the full matrix algebra and every active
sector provides a nonzero matrix. Nondegeneracy and cyclicity of the matrix
trace pairing close each nonzero edge $k\to h$ through a third sector $j$,
producing $h\to j\to k$.\footnote{The formal declarations are
\leanid{MPOTensor.PhysicalSectorFactorization.sectorVirtualMatrix},
\leanid{MPOTensor.PhysicalSectorFactorization.exists_two_edge_return_of_neighboringOperator_ne_zero},
and
\leanid{MPOTensor.PhysicalSectorFactorization.neighboringOperator_reaches_of_ne_zero}.}
Injectivity supplies the spanning hypothesis for every physical-sector
factorization.\footnote{This is
\leanid{MPOTensor.PhysicalSectorFactorization.sectorVirtualMatrixFamily_span_eq_top}.}
Positive-weight endpoint nonvanishing follows directly from the three-site
closure and the Hayashi state decomposition.\footnote{The formal declaration
is \leanid{MPOTensor.exists_active_sectorVirtualMatrix_ne_zero}.}
Therefore the zero-weight-reparameterized inverse-map support is recurrent.
\footnote{This is
\leanid{MPOTensor.zeroWeightReparameterizedInverseMapPhysicalSectorFactorization_isRecurrentSupport}.}

\subsection{Recurrence and per-cycle rescaling}

The sector graph has an edge $k\to h$ exactly when $\eta_{k,h}\neq 0$.
Reachability follows nonzero directed edges, and recurrence means that every
nonzero edge lies on a nonzero directed cycle.\footnote{The formal declarations
are \leanid{MPOTensor.IsSectorEdge},
\leanid{MPOTensor.SectorReaches}, and
\leanid{MPOTensor.IsRecurrentSupport}, in
\texttt{TNLean/MPS/MPDO/SectorRecurrence.lean}.}
The graph is defined from $\eta$, not from its trace matrix $T$. This choice
captures the phase obstruction in
Eqs.~\ref{eq:cpgsv17_mpdo_sal_zcl_eta_local_structure_phase_obstruction_real_edges}
and
\ref{eq:cpgsv17_mpdo_sal_zcl_eta_local_structure_phase_obstruction_imaginary_edge},
which is visible at the level of the neighboring operators.

Reachability by reflexive transitive closure is equivalent to the existence of
a finite directed walk. Recurrence is therefore exactly the return-walk
hypothesis used by the closed-walk coboundary theorem.\footnote{The formal
declarations are
\leanid{TNLean.Algebra.DirectedWalk.reaches_iff_reflTransGen},
\leanid{MPOTensor.sectorReaches_iff_directedWalkReaches}, and
\leanid{MPOTensor.isRecurrentSupport_iff_directedWalkReturns}.}

Every nonempty closed directed walk can be represented by the finite cyclic
family of its source vertices. Consecutive vertices are connected by the
corresponding edges, including the final return to the initial vertex, and the
cyclic product of edge weights equals the walk weight.\footnote{The formal
declarations are
\leanid{TNLean.Algebra.DirectedWalk.closedConsVertices},
\leanid{TNLean.Algebra.DirectedWalk.closedConsVertices_edge}, and
\leanid{TNLean.Algebra.DirectedWalk.weight_closedCons_eq_fin_prod}.}

The required finite-factor rescaling is a general matrix lemma. If a finite
Kronecker product
\begin{align}
    A_0\otimes\cdots\otimes A_{N-1}
        &\geq 0,
    \label{eq:cpgsv17_mpdo_sal_zcl_eta_local_structure_positive_finite_kronecker}
\end{align}
has only nonzero factors, then there are nonzero scalars
$c_0,\ldots,c_{N-1}$ such that
\begin{align}
    \prod_{n=0}^{N-1}c_n
        &=1,
    \label{eq:cpgsv17_mpdo_sal_zcl_eta_local_structure_rescaling_product_one}\\
    c_nA_n
        &\geq 0
        \qquad\text{for every }n.
    \label{eq:cpgsv17_mpdo_sal_zcl_eta_local_structure_rescaled_factors_positive}
\end{align}
The proof removes one factor at a time using the two-factor rescaling lemma.
A companion uniqueness lemma states that if $M\neq 0$ and both $c_1M$ and
$c_2M$ are positive semidefinite, then $c_1/c_2$ is a positive real number.
\footnote{The formal declarations are
\leanid{Matrix.exists_pi_smul_posSemidef_of_finKronecker_posSemidef} and
\leanid{Matrix.exists_pos_real_smul_eq_of_smul_posSemidef}, in
\texttt{TNLean/Algebra/KroneckerFactorPositivity.lean}.}

For a cyclic sector word $(k_n)_n$, define
\begin{align}
    \Omega_k
        &=
        \bigotimes_n\eta_{k_n,k_{n+1}}.
    \label{eq:cpgsv17_mpdo_sal_zcl_eta_local_structure_cyclic_neighbor_product}
\end{align}
Reindexing along the shared horizontal-bond indices identifies $\Omega_k$
entrywise with the finite Kronecker product of the consecutive neighboring
operators around the cycle.\footnote{The formal declaration is
\leanid{MPOTensor.cyclicEtaTensorProduct_submatrix_eq_finKronecker}.}
Combining the finite-factor lemma with positivity of $\Omega_k$ at every
length, formalized by
\leanid{MPOTensor.cyclicEtaTensorProduct_posSemidef}, gives a coherent choice
on any one fixed nonzero cycle: there are nonzero scalars $c_n$ such that
\begin{align}
    \prod_n c_n
        &=1,
    \label{eq:cpgsv17_mpdo_sal_zcl_eta_local_structure_cycle_scalar_product}\\
    c_n\eta_{k_n,k_{n+1}}
        &\geq 0
        \qquad\text{for every }n.
    \label{eq:cpgsv17_mpdo_sal_zcl_eta_local_structure_cycle_edges_positive}
\end{align}
This holds both for an abstract neighboring-operator family and for the
inverse-map-constructed \leanid{MPOTensor.sectorEta}.\footnote{The formal
declarations are
\leanid{MPOTensor.exists_pi_smul_posSemidef_of_cyclicEtaTensorProduct_posSemidef}
and
\leanid{MPOTensor.exists_pi_smul_posSemidef_of_sectorEta_cycle}.}
The result concerns one cycle and does not yet prove that choices from
different cycles agree on a shared edge.

\subsection{The one-sided directed-cut obstruction}

The local-orthogonality contradiction at lines 1465--1470 of Appendix C.2
uses
\begin{align}
    \mathcal A
        &=
        \mathcal A_1+\mathcal A_2,
    \label{eq:cpgsv17_mpdo_sal_zcl_eta_local_structure_subspace_decomposition}\\
    \mathcal A_1\mathcal A_2
        &=0.
    \label{eq:cpgsv17_mpdo_sal_zcl_eta_local_structure_one_sided_product_zero}
\end{align}
It does not require $\mathcal A_2\mathcal A_1=0$
\cite{Cirac2016MPDO_arXiv}. The underlying matrix-algebra statement is
formalized as \leanid{Matrix.submodule_sup_ne_top_of_mul_eq_zero}: two nonzero
linear subspaces of a full complex matrix algebra cannot jointly span the
algebra if every product vanishes in one order.

The zero-weight reparameterization retains every physical summand, sets its
right sector tensor to zero when its Hayashi weight vanishes, and leaves all
physical slices unchanged. For the concrete inverse-map factorization,
positive-weight sectors contain nonzero virtual matrices, injectivity makes the
full sector family span $M_D(\mathbb C)$, and the trace-pairing argument closes
every nonzero edge through a third sector. Recurrence is therefore a
conclusion, not an additional hypothesis.

For the stronger one-component assertion used to prove trace-matrix
primitivity, let $S$ be a nonempty proper forward-closed set of positive-weight
sectors and let $C$ be its complement. The source obtains
$\mathcal A_S\mathcal A_C=0$ from the absence of edges from $S$ to $C$
\cite[Appendix~C.2]{Cirac2016MPDO_arXiv}. This implication is formalized by
\leanid{MPOTensor.PhysicalSectorFactorization.oneSiteSectorSpace_mul_eq_zero}.
The one-site matrices in the directed-cut space are identified with the sector
virtual matrices used by the spanning theorem. On the nonzero-weight subtype,
the spaces associated with $S$ and $C$ are both nonzero, their sum is the full
matrix algebra, and their products in the order
$\mathcal A_S\mathcal A_C$ vanish. The one-sided obstruction proves that the
nonzero neighboring support is strongly connected; no absence of edges from
$C$ to $S$ is needed. Positivity identifies this support with the positive
entries of the restricted real trace matrix, which is therefore irreducible.
The relevant declarations are
\leanid{MPOTensor.PhysicalSectorFactorization.oneSiteSectorMatrix_eq_sectorVirtualMatrix},
\leanid{MPOTensor.PhysicalSectorFactorization.activeSector_reflTransGen_neighboringOperator_ne_zero},
and
\leanid{MPOTensor.PhysicalSectorFactorization.activeSectorTraceMatrix_isIrreducible}.

The periodicity argument differs from the source. At lines 1457--1462,
Ref.~\cite{Cirac2016MPDO_arXiv} passes to a power $T^n$, uses blocking to
remove periodic components, and reduces to primitive diagonal components. The
formal proof does not implement that reduction. Instead, nondegeneracy and
cyclicity of the trace pairing give a closed support walk of length two through
every positive-weight sector. Applying triangle closure to an outgoing edge
gives a closed support walk of length three through the same sector. These
coprime return lengths replace the blocking step.

The length-two statement and positivity of the corresponding diagonal entry
of $T^2$ are formalized by
\leanid{MPOTensor.PhysicalSectorFactorization.exists_active_two_cycle} and
\leanid{MPOTensor.PhysicalSectorFactorization.activeSectorTraceMatrix_pow_two_diag_pos}.
The length-three return is
\leanid{MPOTensor.PhysicalSectorFactorization.activeSectorTraceMatrix_pow_three_diag_pos}.
The theorem
\leanid{Matrix.IsIrreducible.isPrimitive_of_pow_two_diagonal_pos_of_pow_three_diagonal_pos}
combines irreducibility with these two coprime returns and gives a common
exponent $N$ such that
\begin{align}
    (T^N)_{k,h}
        &>0
        \qquad\text{for all active }k,h.
    \label{eq:cpgsv17_mpdo_sal_zcl_eta_local_structure_trace_matrix_primitive_power}
\end{align}
Thus
\leanid{MPOTensor.PhysicalSectorFactorization.activeSectorTraceMatrix_isPrimitive}
completes the active-sector primitivity proof. This replacement of the source's
blocking argument is a local correction, not an additional hypothesis.

\subsection{Coherent vertex rephasing}

Assume recurrence. Fix a nonzero edge $k\to h$. By the uniqueness lemma, every
scalar that makes $\eta_{k,h}$ positive semidefinite agrees with every other
such scalar up to a positive real factor, irrespective of the cycle used to
construct it. For any nonzero cycle, compare the edgewise scalars with the
finite-factor rescaling of that cycle. Their product around the cycle is a
positive real number. The associated phase assignment therefore has trivial
holonomy.

Trivial holonomy makes the edge phases a graph coboundary. Fix a base sector in
each recurrent component and define the phase of a vertex by multiplying edge
phases along a directed path from the base. A return path supplied by
recurrence closes any two candidate paths into a nonzero cycle, so the
definition is path-independent.

The graph-theoretic statement is formalized for a recurrent directed relation:
any group-valued edge assignment whose product is one around every closed
directed walk is a vertex coboundary.\footnote{The formal declaration is
\leanid{TNLean.Algebra.DirectedWalk.exists_vertex_of_closedWalk_weight_eq_one}.}
The proof chooses one representative in each reachability component and
defines vertex weights by path multiplication; a common return path proves
independence of the selected path.

The MPDO-specific bridge is formalized under recurrence. A return path closes
each nonzero neighboring operator into a nonzero cycle. Positive rescaling of
the associated cyclic tensor product supplies a unit phase that makes the edge
operator positive semidefinite. Applying the uniqueness lemma to arbitrary
closed walks shows that the selected edge phases multiply to one around each
walk. The coboundary theorem then gives unit vertex phases $z_k$ satisfying
\begin{align}
    z_k^{-1}z_h\,\eta_{k,h}
        &\geq 0
        \qquad\text{for every }k,h.
    \label{eq:cpgsv17_mpdo_sal_zcl_eta_local_structure_vertex_rephased_eta_positive}
\end{align}
The declarations
\leanid{MPOTensor.exists_vertexPhase_smul_posSemidef} and
\leanid{MPOTensor.exists_rephased_inverseMapPhysicalSectorFactorization}
implement the argument, including dependent sector bond spaces and reciprocal
rephasing of the physical-sector tensors. The latter leaves every complete
cyclic neighboring product unchanged by
\leanid{MPOTensor.PhysicalSectorFactorization.rephase_cyclicNeighboringProduct}.
Combining the positive rephased factorization with the conditional bond
construction gives
\leanid{MPOTensor.nonempty_etaLocalStructureData_of_recurrentSupport}.

This general theorem is explicitly scope-restricted because recurrence is not
a hypothesis of Lemma C.4 of Ref.~\cite{Cirac2016MPDO_arXiv}. The restriction
is eliminated for the concrete inverse-map factorization. The declarations
\leanid{MPOTensor.zeroWeightReparameterizedInverseMapPhysicalSectorFactorization_isRecurrentSupport},
\leanid{MPOTensor.exists_positive_inverseMapPhysicalSectorFactorization}, and
\leanid{MPOTensor.exists_positive_physicalSectorFactorization_of_isSAL}
retain and reparameterize zero-weight sectors, derive recurrence from the
source hypotheses, and construct a coherently positive factorization.

The canonical family
\leanid{MPOTensor.ExplicitEtaOperators.ofHayashiMarkov} does not prove this
result. Its operators are partial traces of normalized sector states, so its
trace matrix is identically one. By contrast, the source inverse-map
construction has
\begin{align}
    T_{k,h}
        &=
        \tr(\eta_{k,h}),
    \label{eq:cpgsv17_mpdo_sal_zcl_eta_local_structure_inverse_map_trace_matrix}
\end{align}
which is generally nonconstant and whose positive-weight restriction is
primitive. The two families cannot be identified without an additional
theorem and are unequal in general.

\subsection{Commuting physical bonds and exact realization}

The translation-invariant bond $B_{n,n+1}$ is defined in the original physical
coordinates and proved positive. Adjacent translates commute on every
periodic chain of length at least three, including the pair crossing the
periodic boundary. Length two requires a separate ordering check because the
translates beginning at sites zero and one read the sites in opposite cyclic
orders. On the $(k,h)$ sector block, they act as
\begin{align}
    B_{0,1}\big|_{(k,h)}
        &=
        I_{L_k\otimes R_h}\otimes\eta_{k,h},
    \label{eq:cpgsv17_mpdo_sal_zcl_eta_local_structure_two_site_forward_bond}\\
    B_{1,0}\big|_{(k,h)}
        &=
        \eta_{h,k}^{\mathrm{op}}
        \otimes I_{R_k\otimes L_h}.
    \label{eq:cpgsv17_mpdo_sal_zcl_eta_local_structure_two_site_reverse_bond}
\end{align}
These operators commute because they act on complementary tensor factors.
Proposition C.8 of Ref.~\cite{Cirac2016MPDO_arXiv} does not discuss this
crossed finite-size ordering separately, although it is an instance of the
translated-bond conclusion.\footnote{The formal declarations are
\leanid{MPOTensor.PhysicalSectorFactorization.physicalBond_adjacent_comm} and
\leanid{MPOTensor.PhysicalSectorFactorization.physicalBond_two_zero_one_comm}.}

Translates with disjoint cyclic two-site supports commute by locality, without
additional physical hypotheses.\footnote{The formal declaration is
\leanid{MPOTensor.PhysicalSectorFactorization.physicalBond_disjoint_comm}.}
The adjacent, crossed two-site, and disjoint cases combine to prove pairwise
commutativity of all translates on every periodic chain of length $N\geq 2$.
\footnote{The formal declaration is
\leanid{MPOTensor.PhysicalSectorFactorization.physicalBond_pairwise_comm}.}

The cyclic sector identity is identified with the ordered bond product and
transported through the local unitary. This gives the exact finite-chain
realization
\begin{align}
    \sigma^{(N)}(K)
        &=
        c_N\prod_{n=1}^{N}B_{n,n+1},
    \label{eq:cpgsv17_mpdo_sal_zcl_eta_local_structure_exact_bond_realization}\\
    c_N
        &=1.
    \label{eq:cpgsv17_mpdo_sal_zcl_eta_local_structure_canonical_bond_scalar}
\end{align}
Here the canonical bond is determined by the physical-sector factorization.
\footnote{The formal declaration is
\leanid{MPOTensor.PhysicalSectorFactorization.mpo_eq_product_physicalBond}.}
Together with positivity of the neighboring operators, this gives the data
required to construct \leanid{MPOTensor.EtaLocalStructureData}. The
conditional construction is
\leanid{MPOTensor.PhysicalSectorFactorization.etaLocalStructureData}, and its
normalization scalar is one. The raw inverse-map physical-sector factorization
and its coherently positive normalization are constructed for every injective
tensor satisfying SAL; ZCL is not needed.

Proposition \texttt{3to4} assumes SAL, not ZCL
\cite{Cirac2016MPDO_arXiv}. ZCL enters only after the commuting product form,
in the later derivation of the GSNNCH--ZCL and renormalization fixed-point
conclusions. Adding ZCL to the conditional theorem would therefore not address
the coherent positive choice.

The coherent positive choice and active-sector trace-matrix primitivity in
Lemma \texttt{propSN} are complete. The latter is formalized by
\leanid{MPOTensor.exists_positive_physicalSectorFactorization_activeSectorTraceMatrix_isPrimitive_of_isSAL}.
The positive commuting bond, its exact all-length product realization, and the
eta-local structure asserted by Proposition \texttt{3to4} follow from the
coherently positive physical-sector factorization. Their composition with the
SAL theorem is
\leanid{MPOTensor.nonempty_etaLocalStructureData_of_isSAL}. Thus the eta-local
conclusion of Proposition C.8 of Ref.~\cite{Cirac2016MPDO_arXiv} is formalized
from its source hypotheses.

The physical sector decomposition and the effective trace-matrix index set
must be distinguished. The physical decomposition retains every summand of the
middle Hilbert space, including those with $p_k=0$. After zero-weight
reparameterization,
\begin{align}
    \eta_{k,h}
        &=0,
        \qquad p_k=0,
    \label{eq:cpgsv17_mpdo_sal_zcl_eta_local_structure_inactive_outgoing_eta_zero}\\
    \eta_{h,k}
        &=0,
        \qquad p_k=0,
    \label{eq:cpgsv17_mpdo_sal_zcl_eta_local_structure_inactive_incoming_eta_zero}
\end{align}
for every $h$. The trace matrix on the ambient sector set therefore has a zero
row and zero column and cannot be primitive. The primitive matrix in Lemma
\texttt{propSN} must be read on the effective sector set
\begin{align}
    \{k:p_k>0\}.
    \label{eq:cpgsv17_mpdo_sal_zcl_eta_local_structure_active_sector_set}
\end{align}
This is the customary convention when zero-probability terms are omitted from
a direct-sum state decomposition. The Lean definition
\leanid{PhysicalSectorFactorization.activeSectorTraceMatrix} restricts the
trace matrix to this subtype while preserving the full physical sector
equivalence and hence the complete Hilbert space.

\section{Classification}

This note records a resolved coherent-positive-choice and primitivity gap in
the SAL-to-positive-bond construction of
\leanid{MPOTensor.EtaLocalStructureData}. Zero-weight sectors are
reparameterized without changing the physical decomposition. Injectivity then
implies recurrence, and a coherent vertex rephasing makes every neighboring
operator positive semidefinite. The general recurrence theorem remains useful
as a scope-restricted graph-theoretic component, but recurrence is a conclusion
for the reparameterized inverse-map factorization.

The trace matrix restricted to positive-weight sectors is primitive. The
ambient trace matrix is not primitive when a zero-weight sector is present.
This distinction preserves all physical summands while identifying the
effective matrix to which the source's primitivity claim applies.

For the conditional bond construction, every pair of translates commutes at
each length $N\geq 2$. The finite-chain MPO is the ordered product of these
bonds, first in cyclic sector coordinates and then in the original physical
coordinates. Positive semidefiniteness of the neighboring operators supplies a
positive commuting bond and an exact realization with scalar one, thereby
determining the eta-local structure data.

An injective tensor satisfying SAL already has the raw physical-sector
factorization. Normalizing inactive sectors and deriving recurrence from
injectivity now provide the coherent positive choice, and the active-sector
trace-matrix primitivity proof is complete. The projector presentation of the
source's local-orthogonality contraction remains a separate presentation, but
it no longer obstructs Lemma \texttt{propSN}: the directed-cut product identity
is proved directly for the sector virtual matrix spaces.

\bibliographystyle{alpha}
\bibliography{references}

\end{document}
